1. 周期引理

字符串的周期 ： $S$ 有周期  $d$ 当且仅当 $\forall i<|S|, S[i] = S[i \bmod d]$ 成立

字符串的完整周期 ： $S$ 有完整周期 $d$ 当且仅当 $d$ 是 $S$ 的周期且满足 $d \mid \lvert S \rvert$

字符串 $S$ 有完整周期 $p$ ， $q$， 那么有完整周期 $\gcd(p,q)$

**Fine-Wilf 定理:**

字符串 $S$ 有周期 $p$ ， $q$，且 $|S| \geq p + q - \gcd(p,q)$ 那么有周期 $\gcd(p,q)$

一个简单的证明：

周期 $p$ 将 $S$ 划分成 $p$ 个等价类， 还要添加 $p - \gcd(p,q)$个等价关系使得 $S$ 可以划分为 $\gcd(p,q)$ 个等价类，而这可以通过在周期 $q$ 后添加  $p - \gcd(p,q)$个字符实现，故需要  $|S| \geq p + q - \gcd(p,q)$ 

边界实例：斐波那契字符串

定义 $f_1 = a , f_2 = b, f_n = f_{n - 1} + f_{n - 2}$ (注意连接顺序)

那么 $g_n = f_n[1: |f_n| - 2]$ 即为边界的实例 （如 $g_7 = babbababbab$,有周期5，8）

2. 判断无穷字符串相等：

如果字符串 $a^{\infty} = b^{\infty}$,那么设字符串 $c$ 是字符串 $a$  和 $b$ 的本原字符串，那么 $a = c^{k_1}, b = c^{k_2}$, 所以我们只用求出来 $a$ 和 $b$ 的最小周期字符串来比较即可。

也可以用下面所描述的算法。

3. 无穷字符串匹配：

如果 $a^{\infty} = b^{\infty}$ ,那么就完全匹配，否则我们可以找第一个失配点，即 $\operatorname{lcp}(a^{\infty},b^{\infty})$的长度。

根据 Fine-Wilf定理，我们只需要检查 $a^{\infty} , b^{\infty}$ 的前 $|a| + |b| - \gcd(|a|,|b|)$个字符。

如果相等，根据 Fine-Wilf定理，我们可以得到 $a^{\infty} , b^{\infty}$的一个共同的周期 $c = \gcd(|a|,|b|)$，那么 $a^{\infty} , b^{\infty}$ 其实都可以表示成 $c^{\infty}$ ,即 $a^{\infty} = b^{\infty}$ 。

否则，我们通过暴力匹配找到了第一个失配点，而暴力匹配的次数不会超过 $|a| + |b| - \gcd(|a|,|b|)$

对于多个无穷字符串 $s^{\infty}_i$ 的匹配/建立字典树问题，我们可以把它们先转化成本原字符串，然后按照字符串长度从大到小插入到trie，但是为了保证出现失配点但不浪费，我们插入的应该是 $s^{\infty}_i [1 : \max(2|s_i|, maxlen)]$ ,其中 $maxlen$ 是字典树中**当前子树**最深节点的深度，所以插入的深度实际上是动态决定的，而至少为长度的2倍则是为了让第一个失配点必须出现（ $2|s_i| \geq |s_i| + |s_j| - \gcd(|s_i|,|s_j|)$）。这保证了我们插入操作是正确的且是最少的。

这里我们给出一个对应的字典树代码